\section{IND-CPA with Decryption Game Details}
\label{app:ind-cpad-game}

This appendix records the game syntax suppressed in
\Cref{sec:games}.
It is not meant to be a second proof outline; the main text reasons about the
named Coq games and their hybrid relations.

\begin{algorithm}
\caption{IND-CPAD Challenger Skeleton}
\label{alg:ind-cpad-game}
Ingredients:
\begin{itemize}
  \item An approximate homomorphic encryption scheme
    $S=(KeyGen,Enc,Eval,Dec)$.
  \item A decrypt-query bound $q\in\bbN$.
\end{itemize}
Global challenger state:
\begin{itemize}
  \item keys $pk,evk,sk$;
  \item hidden challenge bit $b\in\{0,1\}$;
  \item table $T\in(M\times M\times C)^*$ of challenge plaintext pairs and
    produced ciphertexts;
  \item decrypt counter $d$.
\end{itemize}
\begin{algorithmic}
  \Procedure{Game}{$\cA$}
    \State $(pk,evk,sk) \gets KeyGen()$
    \State $b \gets \{0,1\}$
    \State $T \gets [\,]$
    \State $d \gets 0$
    \State $b^* \gets \cA^{OEnc,OEval,ODec}(pk,evk)$
    \State \Return $(b^*=b)$
  \EndProcedure
\end{algorithmic}
\end{algorithm}

\begin{algorithm}
\caption{IND-CPAD Oracles}
\label{alg:ind-cpad-oracles}
The oracles share the challenger state from \Cref{alg:ind-cpad-game}.
\begin{algorithmic}
  \Procedure{OEnc}{$m_0,m_1$}
    \State $c \gets Enc(pk,m_b)$
    \State $T \gets T \doubleplus [(m_0,m_1,c)]$
    \State \Return $c$
  \EndProcedure
  \State
  \Procedure{OEval}{$f,\mathcal{I}$}
    \State $\vec r \gets (T_i)_{i\in\mathcal{I}}$
    \State $m_0' \gets f((r.m_0)_{r\in\vec r})$
    \State $m_1' \gets f((r.m_1)_{r\in\vec r})$
    \State $c' \gets Eval(f,(r.c)_{r\in\vec r})$
    \State $T \gets T \doubleplus [(m_0',m_1',c')]$
    \State \Return $c'$
  \EndProcedure
  \State
  \Procedure{ODec}{$i$}
    \State \textbf{assert} $d<q$
    \State $d \gets d+1$
    \State \textbf{assert} $i<|T|$
    \If{$T_i.m_0\ne T_i.m_1$}
      \State \Return $\mathsf{None}$
    \Else
      \State \Return $\mathsf{SampleCurrent}(T_i)$
    \EndIf
  \EndProcedure
\end{algorithmic}
\end{algorithm}

Here \(\mathsf{SampleCurrent}\) is instantiated by the surrounding game.
The real game returns a noise-flooded sample centered at
$Dec(sk,T_i.c)$.
The simulated-decrypt game returns a noise-flooded sample centered at the
stored plaintext $T_i.m_b$.
The formal proof factors the deterministic prefix of \textsc{ODec} from this
last noisy continuation, which is why the one-call replacement theorem can
focus only on the two Gaussian centers.

\begin{algorithm}
\caption{IND-CPA reduction adversary $\mathcal{B}_{\cA}$}
\label{alg:ind-cpa-reduction}
This is the conventional reduction endpoint used in the main proof.
It stores the same logical table as the IND-CPAD challenger, but obtains
challenge encryptions from the outer IND-CPA encryption oracle.
\begin{algorithmic}[1]
  \Procedure{$\mathsf{Flood}$}{$m,e$}
    \State $\eta \gets D_{\bbZ^n,(e^2+1)\gamma}$
    \State \Return $I_m^{-1}(\eta+I_m(m))$
  \EndProcedure
  \Statex
  \Procedure{$\mathcal{B}_{\cA}$}{$pk,evk$}
    \State $T \gets [\,]$
    \State $d \gets 0$
    \State $b^* \gets
      \cA^{OEnc_{\mathcal{B}},OEval_{\mathcal{B}},ODec_{\mathcal{B}}}(pk,evk)$
    \State \Return $b^*$
  \EndProcedure
  \Statex
  \Procedure{$OEnc_{\mathcal{B}}$}{$m_0,m_1$}
    \State $c \gets OEnc_b(m_0,m_1)$
      \Comment{outer IND-CPA encryption oracle}
    \State $T \gets T \doubleplus [(m_0,m_1,c)]$
    \State \Return $c$
  \EndProcedure
  \Statex
  \Procedure{$OEval_{\mathcal{B}}$}{$f,\mathcal{I}$}
    \State $\vec r \gets (T_i)_{i\in\mathcal{I}}$
    \State $m_0' \gets f((r.m_0)_{r\in\vec r})$
    \State $m_1' \gets f((r.m_1)_{r\in\vec r})$
    \State $c' \gets Eval(f,(r.c)_{r\in\vec r})$
    \State $T \gets T \doubleplus [(m_0',m_1',c')]$
    \State \Return $c'$
  \EndProcedure
  \Statex
  \Procedure{$ODec_{\mathcal{B}}$}{$i$}
    \State \textbf{assert} $d<q$
    \State $d \gets d+1$
    \State \textbf{assert} $i<|T|$
    \State $(m_0,m_1,c) \gets T_i$
    \If{$m_0\ne m_1$}
      \State \Return $\mathsf{None}$
    \EndIf
    \State \textbf{assert} $c=\mathsf{Some}(\_,e)$
    \State $y \gets \mathsf{Flood}(m_0,e)$
    \State \Return $\mathsf{Some}(y)$
  \EndProcedure
\end{algorithmic}
\end{algorithm}

The reduction does not know the IND-CPA challenge bit.
This is harmless for decryption queries that return a value: the simulator
only returns a flooded plaintext when the stored labels agree, $m_0=m_1$, so
there is no bit-dependent plaintext choice left to make.
